Alzheimer's-disease biomarker models are accurate and rarely adopted, because a number without a
rationale cannot justify a lumbar puncture or an anti-amyloid therapy. Chain-of-thought (CoT)
prompting promises an \emph{ante-hoc} rationale emitted by the same forward pass that emits the
prediction, and is widely assumed to make clinical predictions interpretable. We test that
assumption directly. On three Alzheimer's tasks built from open cohorts---neuropathological tau
staging with real quantified amyloid and tau assays, brain-age regression, and longitudinal
dementia progression---we compare CoT against tabular black boxes and against two matched
no-reasoning LLM controls, and we measure whether the reasoning chains causally drive the answers
they accompany. CoT is substantially less accurate than ordinary tabular baselines
($\Delta$AUC $-0.16$ to $-0.18$; brain-age MAE 11.4\,y versus 5.3\,y, worse than predicting the
cohort mean), and on the primary task its chains are not load-bearing: forcing an answer after
showing none of the chain reproduces the full-chain answer 82.9\% of the time, and multiplying a
biomarker value inside the chain by ten changes the answer 3.6\% of the time. A blinded rubric
nonetheless rates those same chains 3.67/4 for clinical plausibility. Three findings complicate a
blanket dismissal: faithfulness is 5.7$\times$ higher on the continuous target than on
classification, cited features shift predictions 1.4--2.1$\times$ more than uncited ones, and CoT
both helps and is faithful at 1.5B parameters while hurting and being decorative at 7B---inverting
the deployment incentive.
